% Nouveau EGA V (cote25): archive cote25, site118 (PDF page119).
% English translation of the checked French annotated master.
% AI-translated by GPT-6 Astra (OpenAI, ChatGPT Pro mode).
\scanpage{25/118}
\mspage{\hand{21}}

\var{\margin{$|$}}{\placedmargin{$|$}}{marginal note placed}{ega25_25_118-v0001}\var{\margin{$\wr$}}{\placedmargin{$\wr$}}{marginal note placed}{ega25_25_118-v0002}which will imply the desired conclusion. Let $\underline{J}$ be the Ideal defining
D; it will obviously suffice to prove that $\underline{J}\hand{\to} i_{\hand{*}}(\underline{J}\hand{|}X-Z")$ is an isomorphism
($i:X-Z"\hand{\to} X$ denoting the canonical immersion); then the homomorphism
$\underline{J}\hand{\to}\underline{O}_X$ is indeed reconstructed by applying the functor $i_{\hand{\unsure{*|\relbar}}}$ to
$\underline{J}\,\hand{|}\,X-Z" \hand{\unsure{\longrightarrow|\relbar}} \underline{O}_X\,\hand{|}\,X-Z"$ . Now since $\underline{J}$ is invertible, it is flat over S \var{\over{and}{\unsure{and}}}{\unsure{and}}{writing over}{ega25_25_118-v0004}
\var{\xout{fibers of}}{}{typewriter strike-out}{ega25_25_118-v0005} $x\hand{\in} Z" \hand{\Rightarrow} \mathrm{prof}_x\,\underline{J}_{s\,\hand{\unsure{*|x}}} \hand{\geqslant} 2$. It therefore suffices to prove the
\uline{Lemma} \hand{3.2}. Let $X\hand{\to} S$ be loc of finite presentation, F a Module of finite
presentation on X flat rel to S, T a closed subset of X; suppose that for
\var{\margin{cf\ \S 20.11}}{\placedmargin{cf\ \S 20.11}}{marginal note placed}{ega25_25_118-v0006}
every $x\hand{\in} X$ over $s\hand{\in} S$, one has \hand{\unsure{/}} $\mathrm{prof}_x F_{s\,\var{\xout{\unsure{n|x}}}{}{typewriter strike-out}{ega25_25_118-v0007}}\hand{\geqslant} 1$ (resp. $\mathrm{prof}_x F_s \hand{\geqslant} 2$).
Then the canonical homomorphism $F\hand{\to} i_{\hand{*}}(F\hand{|}X-T)$ is injective (resp.
bijective), where $i:X-T\hand{\to} T$ is the canonical immersion.

Proof of the lemma: one can suppose S,X affine, and
reduce by the patented procedure to the case of noetherian S. Then the hypothesis
implies, by virtue of par.6 , that $\mathrm{prof}_x F \hand{\geqslant} \mathrm{prof}_x F_s$ for every .
$x\hand{\in} X$ over $s\hand{\in} S$, hence $\mathrm{prof}_x F_s \hand{\geqslant} 1$ (resp.$\hand{\geqslant} 2$) if $x\hand{\in} T$. One then concludes
by par.5 . (NB To do things properly, this lemma should appear in par.
11 under the heading; elimination of noetherian hypotheses ...).

It remains finally to check the last assertion of prop.2, namely
that \var{\deleted{\unsure{Z}}}{}{deletion}{ega25_25_118-v0010}\var{\inserted[typed]{X-Z}}{X-Z}{insertion}{ega25_25_118-v0011} is exactly the \unsure{domain} of definition of the rational
map rel to S defined by f. Let $U\hand{\supset} X-Z$ be this domain; it then follows from
proposition 1 that $U\hand{\longrightarrow} P$ is associated with a\var{\xout{\unsure{family}}}{}{typewriter strike-out}{ega25_25_118-v0013}\var{\inserted[typed]{lin. system}}{ lin. system}{insertion}{ega25_25_118-v0014} of divisors
D' on U\unsure{/|x}S parametrized by $\overset{\hand{\vee}}{P}$ , and one has $D'\hand{|}(X-Z)\times_S\overset{\hand{\vee}}{P} = D\hand{|}(X-Z)\times_S\overset{\hand{\vee}}{P}$ . Applying
the uniqueness result \hand{(}already proved\hand{)} to D' and $D\hand{|}U\times_S\overset{\hand{\vee}}{P}$ , one therefore sees
that these last two divisors are equal, so $D\hand{|}U\times_S\overset{\hand{\vee}}{P}$ has no
fixed points, i.e. $Z\hand{\cap} U=\hand{\emptyset}$, hence U=X-Z , QED.
\tnote{25-118-1}{The overwritten word rendered ``and'' has an uncertain initial e/a in French ``et''. The doubt attached to ``domain'' concerns whether the final e of French ``domaine'' is canceled or remains; the French-only letter variant is recorded as omitted. The doubt in the canceled family phrase concerns the feminine ending e of the preceding French article. The references to prop.2 and proposition1 and the displayed target T are retained. ``Loc'', ``rel'', and ``lin.'' mean locally, relative, and linear. ``Patented procedure'' renders ``procédé bréveté''.}

\mspage{\hand{118}}
